\begin{frame}{Introduction}

  \begin{block}{Sujet du projet}
    Notre projet porte sur \textbf{l'optimisation conjointe des blocs opératoires et des lits d'hospitalisation} (ambulatoire et conventionnelle).

    L'objectif est de concevoir un \textbf{outil d'aide à la décision} capable de planifier les interventions chirurgicales en tenant compte simultanément des ressources du bloc opératoire et des capacités d'hospitalisation.
  \end{block}

  \vspace{0.3em}

  \begin{alertblock}{Pourquoi ce problème ?}
    Le bloc opératoire mobilise des ressources coûteuses et interdépendantes : salle d'opération, chirurgien, équipe d'anesthésie, personnel soignant et lit d'hospitalisation. Une décision prise uniquement à partir du planning du chirurgien peut déplacer la contrainte vers les lits ou les équipes soignantes.
  \end{alertblock}
\end{frame}
\begin{frame}{Introduction}  
  \begin{block}{Objectif global}
    À partir des données historiques de la clinique, notre système :
    \begin{itemize}
      \item prédit la durée d'une intervention et la durée de séjour du patient ;
      \item propose au chirurgien \textbf{deux créneaux optimaux} (une date principale et une alternative) afin d'améliorer la planification hospitalière.
    \end{itemize}
  \end{block}
\end{frame}


% ---------- Flux poussé vs flux tiré ----------
\begin{frame}{Contexte : du flux poussé au flux tiré}

  \begin{columns}[T]

    \begin{column}{0.48\textwidth}

      \begin{block}{Organisation actuelle : flux poussé}
        \begin{itemize}
          \item Le chirurgien choisit une date selon son planning.
          \item Les blocs, les lits et les équipes s'adaptent ensuite.
          \item Cela crée des pics et des creux d'activité.
        \end{itemize}
      \end{block}

    \end{column}

    \begin{column}{0.48\textwidth}

      \begin{block}{Approche proposée : flux tiré}
        \begin{itemize}
          \item La planification dépend des ressources disponibles.
          \item Les blocs et les lits sont optimisés conjointement.
          \item Le chirurgien conserve la décision finale parmi deux propositions.
        \end{itemize}
      \end{block}

    \end{column}

  \end{columns}

  \vspace{0.4em}

  \begin{alertblock}{Idée centrale du projet}
    Passer d'une planification centrée sur le chirurgien à une planification centrée sur la disponibilité réelle des ressources hospitalières.
  \end{alertblock}

\end{frame}